\documentclass[11pt]{article}
\usepackage{hw-solutions-common}
\usepackage{amsmath,amssymb}

\begin{document}

\hwsolhead{Test 2 (PHYS 235, Calculus track)}{Constants used throughout:
$\dfrac{1}{4\pi\varepsilon_0}=8.99\times10^{9}$ N$\cdot$m$^2$/C$^2$,
$\varepsilon_0=8.85\times10^{-12}$ C$^2$/(N$\cdot$m$^2$),
$e=1.602\times10^{-19}$ C, $m_p=1.673\times10^{-27}$ kg.}

\noindent\textbf{Section 1 answer summary:}
\begin{center}
\renewcommand{\arraystretch}{1.3}
\begin{tabular}{c|c|c}
Point & $V$ & $\vec E$ \\\hline
P1 & (C) zero & $+\hat x$ (due right, toward $q$) \\
P2 & (A) positive & $45^\circ$ below $+\hat x$ (southeast) \\
P3 & (A) positive & $0$ \\
\end{tabular}
\end{center}
\vspace{2mm}

%============================================================
\problem{Section 1 --- P1: sign of $V$, direction of $\vec E$}{4}{$Q=+4.0$ nC
sits at the origin, $q=-1.0$ nC sits at $x=5.0$ cm. P1 is on the axis at
$x=4.0$ cm, 4.0 cm from $Q$ and 1.0 cm from $q$.}

\solve{Superpose the two point-charge potentials (dropping the common factor
$1/4\pi\varepsilon_0$ and units, as the hint allows):
\[
V_1 \propto \frac{Q}{r_Q}+\frac{q}{r_q} = \frac{4.0}{4.0}+\frac{-1.0}{1.0} = 1.0-1.0 =
\boxed{0}.
\]
\textbf{(C) zero.} For the field: $Q$ is positive and P1 is to its right, so
$\vec E_Q$ points away from $Q$, in $+\hat x$. $q$ is negative and P1 is to
its left, so $\vec E_q$ points \emph{toward} $q$ -- also $+\hat x$. The two
contributions add instead of fighting each other, so
\[
\boxed{\vec E_1 \text{ points in } +\hat x \text{ (due right, toward } q\text{)}}.
\]}

\review{P1 is the case where $V=0$ but $\vec E\neq 0$ -- the two potentials
cancel exactly (closer to the smaller-magnitude charge by exactly the ratio
that makes it work), but the two field contributions don't fight each other
at all, they add. $V=0$ and $\vec E=0$ are not the same condition, and this
point is designed to make that visible.}

%============================================================
\problem{Section 1 --- P2: sign of $V$, direction of $\vec E$}{4}{P2 is at
$(5.0,\,2.4)$ cm -- directly above $q$, 2.4 cm away from it, and
$\sqrt{5.0^2+2.4^2}=5.5$ cm from $Q$.}

\solve{\[
V_2 \propto \frac{4.0}{5.5}+\frac{-1.0}{2.4} = 0.73-0.42 = \boxed{+0.30}
\;\Rightarrow\; \textbf{(A) positive}.
\]
For the field, weight each unit vector by $1/r^2$. $\vec E_Q$ points away
from $Q$, along the line from $Q$ to P2 (up and to the right); $\vec E_q$
points straight down toward $q$ (P2 sits directly above it). Using
$ E\propto |charge|/r^2$:
\[
\vec E_{Q,2}\propto \frac{4.0}{5.5^2}(0.90,\,0.43) = (0.117,\,0.056),
\qquad
\vec E_{q,2}\propto \frac{1.0}{2.4^2}(0,\,-1) = (0,\,-0.174).
\]
Adding, $\vec E_2 \propto (0.117,\,-0.117)$ -- equal magnitude $+x$ and
$-y$ components, i.e.
\[
\boxed{\vec E_2 \text{ points } 45^\circ \text{ below } +\hat x \text{ (southeast)}}.
\]}

\review{The $2.4$ cm height isn't an arbitrary number -- it's exactly the
value that makes $Q$'s field and $q$'s field add up to a clean $45^\circ$
direction, so a numeric sanity check (rather than just a qualitative
``down-and-to-the-right'' guess) is the fastest way to confirm you have the
right compass direction.}

%============================================================
\problem{Section 1 --- P3: sign of $V$, direction of $\vec E$}{4}{P3 is on
the axis at $x=10$ cm, beyond $q$: 10 cm from $Q$, 5.0 cm from $q$.}

\solve{\[
V_3 \propto \frac{4.0}{10}+\frac{-1.0}{5.0} = 0.40-0.20 = \boxed{+0.20}
\;\Rightarrow\; \textbf{(A) positive}.
\]
For the field: $\vec E_Q \propto Q/r_Q^2 = 4.0/10^2 = 0.040$, pointing away
from $Q$ (i.e.\ $+\hat x$, since P3 is to the right of $Q$). $\vec E_q
\propto |q|/r_q^2 = 1.0/5.0^2 = 0.040$, pointing \emph{toward} $q$ (i.e.\
$-\hat x$, since P3 is to the right of $q$). Equal magnitude, opposite
direction:
\[
\boxed{\vec E_3 = 0}.
\]}

\review{P3 is the mirror image of P1: there, $V$ canceled but $\vec E$
didn't; here $\vec E$ cancels exactly (by design: $4.0/10^2$ and
$1.0/5.0^2$ both equal $0.040$) while $V$ stays comfortably positive. Together
the three points cover all three cases a student needs to tell apart:
$V=0,\vec E\neq0$ (P1); both nonzero (P2); $V\neq0,\vec E=0$ (P3).}

%============================================================
\problem{Section 2 --- Part 1: the dipole moment}{5}{Charges of magnitude
5 nC, separated by $6\ \mu$m, with $\vec E = 20$ N/C $\hat y$. The dipole
moment points $60^\circ$ below the $x$-axis (``Orientation A'').}

\solve{The dipole moment has magnitude
\[
p = qd = (5.0\times10^{-9}\text{ C})(6.0\times10^{-6}\text{ m}) =
\boxed{3.0\times10^{-14}\text{ C}\cdot\text{m}},
\]
pointing (by definition, from $-q$ to $+q$) $60^\circ$ below the $+x$-axis,
i.e.
\[
\vec p = p(\cos(-60^\circ),\,\sin(-60^\circ),\,0) = (1.5,\,-2.6,\,0)\times10^{-14}\text{
C}\cdot\text{m}.
\]
\textbf{Drawing:} the $+x$-axis horizontal, $+y$ vertical. $\vec E$ is a
vertical arrow pointing straight up ($+\hat y$). $\vec p$ is a separate arrow
starting at the dipole's center, $60^\circ$ \emph{below} the $+x$-axis (down
and to the right, into the fourth quadrant) -- \emph{not} aligned with
$\vec E$ and not symmetric about either axis. The negative charge sits at
the tail of $\vec p$, the positive charge at the head, $6\ \mu$m apart along
that same line.}

\review{A common error here is drawing $p$ at $60^\circ$ \emph{above} the
axis, or measuring the angle from $\vec E$ instead of from the $x$-axis --
read the setup's wording carefully (as the problem itself warns): the angle
given is always relative to the axis stated, not to the field.}

%============================================================
\problem{Section 2 --- Part 2: torque}{4}{Same dipole and field as above.}

\solve{$\vec\tau = \vec p\times\vec E$. Since $\vec p=(p_x,p_y,0)$ and
$\vec E=(0,E_y,0)$ (purely $\hat y$),
\[
\vec\tau = (0,\,0,\,p_xE_y) = (0,\,0,\,(1.5\times10^{-14})(20)) =
\boxed{(0,\,0,\,3.0\times10^{-13})\text{ N}\cdot\text{m}},
\]
i.e.\ $3.0\times10^{-13}$ N$\cdot$m in $+\hat z$ (out of the page).}

\review{The torque vector's direction is the rotation \emph{axis}, found by
the right-hand rule, not a push in some in-plane direction. $+\hat z$ means:
curl the right hand's fingers from $\vec p$ toward $\vec E$ through the
\emph{smaller} angle between them, and the thumb points out of the page --
confirming the dipole is being spun counterclockwise (as drawn, with $+x$
right and $+y$ up), rotating $\vec p$ up and to the left, toward alignment
with $\vec E$.}

%============================================================
\problem{Section 2 --- Part 3: potential energy, Orientation A vs.\ B}{6}{Same
dipole and field. Orientation B is whichever orientation has the lowest PE.}

\solve{$U=-\vec p\cdot\vec E=-pE\cos\theta$, where $\theta$ is the angle
between $\vec p$ and $\vec E$. This is minimized when $\cos\theta=1$, i.e.\
$\vec p$ points the \emph{same direction as} $\vec E$ -- straight up,
$+\hat y$. \textbf{Orientation B: $\vec p$ parallel to $\vec E$ (pointing
straight up).}

In Orientation A, $\vec p$ is at $-60^\circ$ and $\vec E$ is at $+90^\circ$
(both measured from $+x$), so $\theta_A = 90^\circ-(-60^\circ)=150^\circ$:
\[
U_A = -pE\cos150^\circ = -(3.0\times10^{-14})(20)(-0.866) = +5.2\times10^{-13}\text{ J}.
\]
In Orientation B, $\theta_B=0$:
\[
U_B = -pE\cos0^\circ = -(3.0\times10^{-14})(20) = -6.0\times10^{-13}\text{ J}.
\]
So Orientation A has
\[
U_A-U_B = 5.2\times10^{-13}-(-6.0\times10^{-13}) = \boxed{1.1\times10^{-12}\text{
J \emph{more} than Orientation B}}.
\]}

\review{$U_A-U_B = pE(1-\cos\theta_A)$ ranges from $0$ (if A already equals
B) to $2pE=1.2\times10^{-12}$ J (if A is the \emph{worst} orientation,
anti-aligned with $\vec E$). $150^\circ$ is close to that worst case
($180^\circ$), so it makes sense that $U_A-U_B$ comes out close to, but
just under, that $2pE$ ceiling.}

%============================================================
\problem{Section 3 --- Part 1: field at A, B, C}{10}{Positive rod
($\lambda_0$) and negative rod ($-4\lambda_0$) separated by $2d$. A is a
distance $d$ to the left of the positive rod; C is a distance $d$ to the
right of the negative rod; B is halfway between (distance $d$ from each).}

\solve{Put the positive rod at $x=0$ and the negative rod at $x=2d$, so
$A=-d$, $B=d$, $C=3d$. All three points lie on the line through both rods,
so each rod's field there is purely along $\hat x$: pointing \emph{away}
from the positive rod, \emph{toward} the negative rod.

\textbf{At A} ($r_{\text{pos}}=d$, $r_{\text{neg}}=3d$): the positive rod's
field pushes left ($-\hat x$, away from it); the negative rod's field pulls
right ($+\hat x$, toward it, since the negative rod is to A's right).
\[
E_A = -\frac{\lambda_0}{2\pi\varepsilon_0 d} + \frac{4\lambda_0}{2\pi\varepsilon_0(3d)}
= \frac{\lambda_0}{2\pi\varepsilon_0 d}\Big(-1+\tfrac43\Big) =
\boxed{\frac{\lambda_0}{6\pi\varepsilon_0 d}\ \ (+\hat x\text{, toward the rods})}.
\]

\textbf{At B} ($r_{\text{pos}}=r_{\text{neg}}=d$): now \emph{both} fields
point the same way -- away from the (weaker) positive rod is $+\hat x$ here
(B is to its right), and toward the (stronger) negative rod is also
$+\hat x$ (B is to its left). They add, not cancel:
\[
E_B = \frac{\lambda_0}{2\pi\varepsilon_0 d} + \frac{4\lambda_0}{2\pi\varepsilon_0 d} =
\boxed{\frac{5\lambda_0}{2\pi\varepsilon_0 d}\ \ (+\hat x)}.
\]

\textbf{At C} ($r_{\text{pos}}=3d$, $r_{\text{neg}}=d$): away from the
positive rod is $+\hat x$ (C is to its right); toward the negative rod is
$-\hat x$ (C is to its right of it, so ``toward'' means back to the left).
\[
E_C = \frac{\lambda_0}{2\pi\varepsilon_0(3d)} - \frac{4\lambda_0}{2\pi\varepsilon_0 d}
= \frac{\lambda_0}{2\pi\varepsilon_0 d}\Big(\tfrac13-4\Big) =
\boxed{-\frac{11\lambda_0}{6\pi\varepsilon_0 d}\ \ \big(\text{i.e.\ } \tfrac{11\lambda_0}{6\pi\varepsilon_0 d}\text{ in } -\hat x\text{, toward the rods}\big)}.
\]
With $d=6$ cm and $\lambda_0=400$ nC/m (the 80\%-credit numeric option):
$E_A \approx 4.0\times10^4$ N/C, $E_B\approx6.0\times10^5$ N/C, $E_C\approx
4.4\times10^5$ N/C.}

\review{B and C both end up dominated by the much closer, 4$\times$-stronger
negative rod (that's why both exceed $E_A$ by more than an order of
magnitude), but A is the one case where the two contributions subtract
instead of reinforcing -- which is exactly why $A$'s field comes out so much
smaller than the others, even though it's also only a distance $d$ from its
near rod, same as $B$ and $C$.}

%============================================================
\problem{Section 3 --- Part 2: is there a location where $\vec E=0$?}{6}{Same
two rods, now fixing $d=6$ cm.}

\solve{First: \emph{where} could a null point even be possible? At any point
off the line through the two rods, the two rods' fields point in genuinely
different (non-opposite) directions -- two nonzero vectors that aren't
anti-parallel can never sum to zero, regardless of their magnitudes. So a
null point, if one exists, must lie on the line through both rods (the same
line $A$, $B$, $C$ sit on).

On that line there are three regions to check. \emph{Between the rods}
($0<x<2d$): both fields point the same way (as at $B$ above) -- they can
only add, never cancel, so no zero there. \emph{Right of the negative rod}
($x>2d$, like $C$): setting the two magnitudes equal,
$\lambda_0/r_{\text{pos}} = 4\lambda_0/r_{\text{neg}}$ with $r_{\text{neg}} =
r_{\text{pos}}-2d$, gives $x=-2d/3$ -- negative, so \emph{not} actually in
this region. No zero here either.

\emph{Left of the positive rod} ($x<0$, like $A$): here the positive rod's
field points further left ($-\hat x$) and the negative rod's field points
right ($+\hat x$, toward it) -- opposite directions, so cancellation is at
least possible. Setting magnitudes equal, with $r_{\text{pos}}=|x|=-x$ and
$r_{\text{neg}}=2d-x$:
\[
\frac{\lambda_0}{-x} = \frac{4\lambda_0}{2d-x}
\;\Rightarrow\;
2d-x = -4x
\;\Rightarrow\;
\boxed{x=-\tfrac{2d}{3}\text{, i.e.\ a distance }\tfrac{2}{3}d\text{ to the
left of the positive rod}}.
\]
With $d=6$ cm, that's $4.0$ cm to the left of the positive (left) rod --
between the positive rod and point $A$ (which sits a full $d=6$ cm out).}

\review{It makes sense the null point sits \emph{closer} to the positive
rod than $A$ does: since the negative rod is $4\times$ stronger, the field
point has to huddle in close to the weaker rod -- well inside $A$'s
distance -- for the weaker rod's $1/r$ advantage (from proximity) to
offset the stronger rod's $4\times$ charge advantage.}

%============================================================
\begin{center}{\LARGE\bfseries Section 4: Gauss's Law}\end{center}
\vspace{2mm}

\problem{Part 1: enclosed charge}{8}{A cube with 3.0 cm sides. The field is
purely in $\hat y$ everywhere. At the top face, $\vec E=\langle
0,-34,0\rangle$ N/C; at the bottom face, $\vec E=\langle 0,+20,0\rangle$
N/C.}

\solve{Since $\vec E$ has no $x$ or $z$ component, the flux through all four
side faces is zero (their outward normals are $\pm\hat x$ or $\pm\hat z$,
perpendicular to $\vec E$). Only the top and bottom faces contribute. With
$A=(0.030\text{ m})^2 = 9.0\times10^{-4}$ m$^2$:

\emph{Top face}, outward normal $+\hat y$:
$\Phi_{\text{top}} = \vec E_{\text{top}}\cdot\hat n\,A = (-34)(+1)A = -34A$.

\emph{Bottom face}, outward normal $-\hat y$:
$\Phi_{\text{bottom}} = \vec E_{\text{bottom}}\cdot\hat n\,A = (+20)(-1)A = -20A$.

\[
\Phi_{\text{total}} = -34A-20A = -54A = -54(9.0\times10^{-4}) = -0.0486\text{
N}\cdot\text{m}^2/\text{C}.
\]
By Gauss's law,
\[
Q_{\text{enc}} = \varepsilon_0\Phi_{\text{total}} =
(8.85\times10^{-12})(-0.0486) = \boxed{-4.3\times10^{-13}\text{ C} \approx
-0.43\text{ pC}}.
\]
A drawing should show the cube with $\vec E$ arrows pointing \emph{down}
into the top face (since $\vec E_{\text{top}}$ points in $-\hat y$, into the
cube from above) and pointing \emph{up} into the bottom face (since
$\vec E_{\text{bottom}}$ points in $+\hat y$, into the cube from below) --
field lines converging into the cube from both the top and the bottom,
consistent with a net negative charge trapped inside.}

\review{The two faces don't need to have equal field strength for this to
work -- Gauss's law only cares about the net flux. Here $34$ and $20$ N/C
aren't equal, so the two faces don't exactly cancel, and that imbalance
(both pointing inward, but by different amounts) is exactly what leaves a
nonzero, negative $Q_{\text{enc}}$.}

\problem{Part 2: what Gauss's law can and can't tell you}{4}{Same cube.}

\solve{Gauss's law only ever gives the \emph{total} charge enclosed by a
surface -- it has no way to say \emph{where} inside the cube that charge
actually sits. The $-0.43$ pC found above could be a thin sheet right at the
cube's center, spread uniformly through the whole volume, concentrated near
one face, or any other distribution; this calculation can't distinguish
between them, and nothing in the given field values (just two numbers, at
the top and bottom) provides enough information to localize it further.

For charge \emph{outside} the cube: this particular Gaussian surface (the
cube itself) tells us nothing about it. To learn anything about charge
outside, you'd need to draw a \emph{different} Gaussian surface that
encloses some of that outside region, together with field information on
that new surface -- neither of which this problem gives us. The one thing
we \emph{can} say is that whatever is outside is consistent with the field
being purely $\hat y$-directed everywhere (not just inside the cube): that's
a strong symmetry statement (reminiscent of infinite charged sheets, where
the field depends only on height, never on the horizontal position), but it
still doesn't pin down a location or an amount.}

\review{This pair of questions is really testing the same idea from two
sides: Gauss's law is a statement about a \emph{surface integral} of the
field, which only ever recovers the enclosed charge's \emph{total}, not its
spatial distribution, and only for the one surface you chose to apply it to.}

\end{document}
